\documentclass{article}

\usepackage[pagebackref=true,colorlinks,urlcolor=magenta]{hyperref}
\usepackage[capitalize,noabbrev]{cleveref}

\usepackage[left=1.25in, top=1in, bottom=1in, right=1.25in]{geometry}
\parskip 2.mm
\parindent 0.mm




\begin{document}


\section*{Cover Letter}

Dear Editors and Reviewers,

\vspace{0.2cm}

\textbf{``Event Camera Vision in the Era of Large Models: A Survey"}

by Lingdong Kong, Haiqian Han, Lai Xing Ng, Xiangyang Ji, Wei Tsang Ooi, and Benoit R. Cottereau

which is submitted to the IEEE Transactions on Pattern Analysis and Machine Intelligence (TPAMI) for consideration for publication.


\subsection*{Summary of Contributions}

Event cameras are bio-inspired neuromorphic sensors that asynchronously capture per-pixel brightness changes with microsecond temporal resolution, high dynamic range, and ultra-low power consumption. While several surveys have reviewed aspects of event-based vision, none adequately addresses the transformative impact of large-scale pre-trained models --- including vision transformers, diffusion models, CLIP, SAM, and multimodal large language models --- on this rapidly evolving field. Our manuscript fills this gap by providing a comprehensive, structured, and up-to-date review.


The key contributions of this work are summarized as follows:
\begin{itemize}
    \item We provide \textbf{comprehensive and timely coverage} of over 400 papers spanning the full spectrum of event camera vision, with particular emphasis on methods from 2024--2026 that leverage foundation models, diffusion priors, state space models, and 3D Gaussian Splatting.

    \item We introduce a \textbf{unified five-pillar taxonomy} that organizes the literature around benchmarks and datasets, perception, reconstruction, understanding, and applications --- providing a coherent framework connecting low-level event processing to high-level semantic reasoning.

    \item We present the \textbf{first systematic review of event camera understanding}, covering event-based MLLMs (EventGPT, EventVL, EventFlash), CLIP-based recognition, open-vocabulary segmentation (OpenESS, SEAL), and SAM/DINO adaptation --- an area not covered by any prior survey.

    \item We identify \textbf{paradigm shifts} driven by large models in each area: from discriminative to generative reconstruction, from supervised to open-vocabulary perception, from NeRF to 3DGS for 3D representation, and from task-specific to foundation-model-based understanding.

    \item We maintain a \textbf{continuously updated companion repository} (\url{https://github.com/worldbench/awesome-event-camera-vision}) with curated paper lists, dataset links, and code repositories.
\end{itemize}

Enclosed, please find the manuscript of this submission.

We kindly ask for your consideration. We are open to making any necessary revisions based on the feedback from the reviewers and ensuring our work aligns with the high standards of TPAMI.

Lastly, we thank the reviewers and editors for the time and effort devoted to this review.

\vspace{0.3cm}

Yours sincerely,

\vspace{0.3cm}

Lingdong Kong, Haiqian Han, Lai Xing Ng, Xiangyang Ji, Wei Tsang Ooi, and Benoit R. Cottereau

\end{document}
